\scanpage{80}\mspage{39}
% archive sheet number: 79
les structures supplémentaires reliant les données $S,F$ etc \unsure{à|:} $\Lambda$ (cf. (13)(14)(5) et aussi (118) p. 23)
% slanted
\margin{il y a aussi l'opération de $\Lambda^*$ \deleted{\unsure{par}} sur $\widetilde S_M$}

Enfin, dans le cas où on ne suppose plus $\Lambda=\Lambda_0$, on voudrait introduire sur $\widetilde P_{M,i}$ \inserted{resp. sur} $\widetilde P_M$ un élément de structure supplémentaire (peut-être l'opération de $\Lambda^*$ sur $\widetilde S_M$) qui permette néanmoins d'interpréter $\mathrm{Gl}_\Lambda(M)^{\pm}\deleted{\unsure{(\ )}}$ resp. $\mathrm{Gl}_\Lambda(M)$ comme le groupe d'automorphismes de cette structure.

\deleted{\unsure{Si $2.1_\Lambda\ne0$}} On suppose dorénavant $2.1_\Lambda\ne0$.
\deleted{B) Le groupe $\{\pm\operatorname{id}_M\}$ opère librement sur $\widetilde S_M,\widetilde A_M,\widetilde F_M$ [\uline{NB} si $\{x,y,z\}\deleted{\unsure{\}}}\in\widetilde F_M$, alors les éléments $x,y,z,-x,-y,-z$ sont tous distincts...]. On pose}

Revenons à $\widetilde P_M$ \inserted{dans le cas $\Lambda$ quelconque}, et considérons l'opération de $\Lambda^*$ sur $\widetilde{\deleted{\unsure{B}}S}_M$, par les $\widetilde T_\lambda$
\[
 \widetilde T_\lambda x=\lambda.x\tag*{(20)}
\]
alors $\Lambda^*$ respecte toutes les structures ($\deleted{\unsure{a}}\widetilde F_M\subset\mathfrak P_3(\widetilde S_M)$ et la partition de $\widetilde F_M$ en parties $\widetilde F_{M,i}$) -- l'opération de $\mathrm{Gl}_\Lambda(M)$ sur $\widetilde S_M$ \deleted{est} commute à $\Lambda^*\inserted{\simeq\operatorname{Centre}(\mathrm{Gl}_\Lambda(M))}$, en respectant aussi
% slanted
\margin{\uline{NB} la donnée de l'antipodisme est contenue dans cette opération de $\Lambda^*$}
